\section{Study design and protocol freeze}
\label{sec:design}

This section is placed before the results deliberately. The design was fixed in
writing before the series ran, and the sequence
\emph{hypothesis $\rightarrow$ frozen method $\rightarrow$ data $\rightarrow$
interpretation} is what allows a reader to distinguish a finding from a story
assembled after the fact.

\subsection{What was frozen, and when}

The pre-registration is a repository artefact dated before the first run of the
series\code. It fixes the question, the hypotheses with their falsification
conditions, the intervention, the dependent variables, the confounders and
their controls, the sample size with its power rationale, the stopping rule,
the statistical tests, and the exclusion policy. Its amendment section is
empty: nothing was changed after the first run started.

Three commitments in it do the real work:

\begin{enumerate}
\item \textbf{The verdict is external.} A run's outcome is the exit status of
      the task's own verifier, never the agent's report of its own success. A
      run without a verifier verdict does not enter the dataset at all.
\item \textbf{Tests were chosen before the data.} Paired exact McNemar for
      capability, exact sign test on per-pair ratios for cost, Spearman rank
      correlation for the horizon trend. No test was selected after seeing an
      effect.
\item \textbf{Exclusions were defined before the counts were seen.} The
      contamination rule (\S\ref{sec:pitfalls}) and the duplicate-collapse rule
      were both written before they were applied. The latter keeps the
      \emph{earliest} launch of a re-run cell, so a repeat cannot promote
      itself into the sample by winning.
\end{enumerate}

\subsection{Hypotheses and how each could fail}

\begin{description}[style=nextline]
\item[H1 (cost).] The median cost ratio of state-only to control exceeds 1.0.
  \emph{Refuted if} the ratio is $\le 1.0$, the sign test misses significance,
  or the direction is inconsistent across the three trials.
\item[H2 (horizon).] The cost ratio falls as horizon grows, crossing 1.0:
  the overhead amortises on long tasks. \emph{Refuted if} the correlation is
  absent, positive, or never approaches 1.0 in the observed range.
\item[H3 (capability).] Reported with its interval and treated as underpowered
  unless the observed discordance exceeds the pre-computed requirement of six
  one-sided discordant pairs.
\end{description}

H2 is the hypothesis this study refutes in the pre-registered manner, and we
report it as such in \S\ref{sec:economics} rather than quietly dropping it.

\subsection{Multiple comparisons}

Three hypotheses were pre-registered, and those three carry the study's
confirmatory weight. Everything else this paper reports a $p$-value for, the
per-stratum splits, the ladder steps of \S\ref{sec:composition}, the
per-condition cost comparisons, is secondary and exploratory, and no
family-wise correction is applied to it.

We state this explicitly because \S\ref{sec:gaps} criticises the original for
not doing so, and the same standard has to apply here. The consequence is
concrete: our headline result survives any correction one would care to apply
($p={}$\McnPTailZero{} on the primary comparison, against a Bonferroni
threshold of $0.05/3$ for the confirmatory family), but several secondary
results sit close enough to $0.05$ that they would not. Where a secondary
result is load-bearing for a claim in this paper we say so and treat it as
provisional rather than established.

\subsection{A pre-declared artefact trap}

One trap was named in advance because it would have produced a flattering
result: if the state-only arm showed a \emph{lower} cost on long tasks merely
because long tasks are where it more often exhausts its turn budget and stops
early, that is censoring, not amortisation. The pre-registration therefore
excludes budget-stopped runs from cost aggregation and reports their count per
stratum separately. The trap was real: the stop-reason distribution in
Table~\ref{tab:stop} shows budget stops concentrated in exactly that arm.

\subsection{Scope}

This series addresses the cost and horizon questions, and reports capability.
Schema sensitivity and state composition are separate designs with their own
pre-registrations; their status is given in \S\ref{sec:composition}.
\pending{if the owner elects to publish E1 alone, this paragraph becomes the
forward reference to the second paper}
